% Generated draft: requires mathematical and source review.
\documentclass[11pt]{article}
\usepackage[margin=25mm]{geometry}
\usepackage{fontspec}
\setmainfont{Latin Modern Roman}
\usepackage{amsmath,amssymb,mathtools,unicode-math}
\setmathfont{Latin Modern Math}
\usepackage{enumitem,needspace,xurl,hyperref}
\setlength{\parindent}{0pt}
\setlength{\parskip}{5pt}
\setlength{\emergencystretch}{3em}
\newcommand{\N}{\mathbb{N}}
\newcommand{\Z}{\mathbb{Z}}
\newcommand{\Q}{\mathbb{Q}}
\newcommand{\R}{\mathbb{R}}
\newcommand{\C}{\mathbb{C}}
\newcommand{\F}{\mathbb{F}}
\newcommand{\A}{\mathbb{A}}
\newcommand{\PP}{\mathbb{P}}
\newcommand{\E}{\mathbb{E}}
\newcommand{\eps}{\varepsilon}
\newcommand{\dd}{\,\mathrm{d}}
\begin{document}
\section*{UnsolvedMath ID: 30001475}
\textbf{Source problem code:} catalogue entry
\subsection{Definitions and mathematical statement}
Use the coordinate model \(\mathbb H^1=(\R^3,\cdot)\) with group law
\[
(x,y,t)\cdot(x',y',t')=
\left(x+x',y+y',t+t'+\frac{xy'-yx'}2\right).
\]
Its horizontal distribution is spanned by
\[
X=\partial_x-\frac y2\partial_t,
\qquad
Y=\partial_y+\frac x2\partial_t,
\]
with \(X,Y\) orthonormal. The associated Carnot--Carath\'eodory distance is defined by minimizing the lengths of horizontal curves. This fixes a coordinate normalization for the sub-Riemannian structure.

For notation, let
\[
W=\{(0,y,t):y,t\in\R\},\qquad
V=\{(x,0,0):x\in\R\}.
\]
Given a domain \(D\subseteq W\) and a function \(u:D\to\R\), the intrinsic graph associated with this splitting is
\[
\Gamma_u=\{w\cdot(u(w),0,0):w\in D\}.
\]
The graphing direction and coordinate normalization here provide notation; the supplied record does not specify the admissible domains or regularity of \(u\).

Sub-Riemannian area refers to horizontal area, or to the corresponding horizontal perimeter when surfaces are treated as boundaries of sets. For a sufficiently regular surface \(\Sigma\), its horizontal area in an open set \(U\) can be written
\[
A_H(\Sigma;U)=\int_{\Sigma\cap U}|N_H|\,\mathrm dA_g,
\]
where \(g\) is the left-invariant Riemannian metric making \(X,Y,\partial_t\) orthonormal, \(N\) is a \(g\)-unit normal, and \(N_H\) is its projection onto the horizontal distribution. For a specified class of variations, area-stationarity means vanishing first variation of horizontal area, and stability means nonnegative second variation. Area minimization requires an area comparison with a specified class of competitors. The source does not specify the regularity, competitor class, or weak formulation needed to apply these notions to every graph under consideration.

\textbf{Requested classification.} Classify all complete area-minimizing intrinsic graphs in the sub-Riemannian Heisenberg group \(\mathbb H^1\). In particular, determine whether the classification referred to in the supplied record for complete stable area-stationary intrinsic graphs extends to all complete area-minimizing intrinsic graphs.

\textbf{Unresolved qualifications.} This is a classification request whose precise target class is incompletely specified by the supplied record. The record does not state the classification to be extended or its hypotheses. It also does not specify whether graphs must be entire, which notion of completeness is intended, whether boundary is allowed, or which regularity and minimization assumptions apply. In particular, completeness for the restricted ambient distance and completeness for a surface length metric must not be silently identified. The displayed definitions provide a conventional framework, but do not resolve these missing hypotheses or assert that any particular classification applies.
\subsection{Short English statement}
The problem asks for a classification of complete intrinsic graphs in \(\mathbb H^1\) that minimize sub-Riemannian area, and whether the classification mentioned for complete stable area-stationary intrinsic graphs applies to them. The supplied record omits the comparison classification and essential hypotheses concerning regularity, completeness, domains, boundary, and competitors, so the intended precise scope requires mathematical review.
\subsection{Sources}
\begin{itemize}
\item Supplied local UnsolvedMath registry record, permanent ID \(30001475\), source problem code \texttt{catalogue entry}, title \emph{Area-Minimizing Intrinsic Graphs in the Heisenberg Group}; collection: Oberwolfach Reports Open Problems. The record gives proposed year 2010 and names no individual proposer.
\item Supplied primary report citation: \emph{Progress in Surface Theory} (2010), Oberwolfach Reports, pp.\ 1225--1318, DOI 10.4171/owr/2010/21. \url{https://doi.org/10.4171/owr/2010/21}. The background describes Oberwolfach Reports as documenting mathematical workshops and their open-problem discussions; no problem-specific page is supplied.
\item Supplied literature source labeled ``OWR 2010/21'': \url{https://doi.org/10.4171/owr/2010/21}.
\item Supplied literature source labeled ``Barone Adesi--Serra Cassano--Vittone'': \url{https://doi.org/10.1007/s00526-006-0076-3}. No title or theorem statement is supplied.
\item Supplied literature source labeled ``Serra Cassano--Vedovato'': \url{https://cvgmt.sns.it/media/doc/paper/4536/SerraCassano_Vedovato_cvgmt_26_11_2019.pdf}. No title or theorem statement is supplied.
\item Literature assessment date: \texttt{2026-08-21}. This is supplied dated curation, not an independently verified assessment. It labels the problem open and states that intrinsic-graph Bernstein and area-minimization results cover important regularity and dimension regimes, while the broad classification question remains unresolved. The background explicitly requests independent verification of time-sensitive status and cited claims.
\item The registry labels the statement \texttt{corrected\_verified} and reports that its mathematical content was checked against the official Oberwolfach report, with notation, context, and question form normalized without changing the source-level claim. This is supplied verification metadata; no source visit or independent verification is asserted here.
\end{itemize}
\end{document}
